\chapter{Data Layer}
\label{cap:datalayer}

The Data Layer, also referred as Data Access Layer (DAL) represent the access and handling to data from one or more source such as different relational database, NoSQL database, file systems or others persistent storage.

Data Layer commonly used to centralize data access logic, reducing redundances and providing a consistent interface for retrieving or storing data.


For GeoMedia, the data layer is implemented using a Relational Database Management System (RDBMS), that represent the most common and well-established approach for the structured management, storage and retrieval of data. The RDBMS organizes information into logically related tables composed of rows and columns, while providing mechanisms for defining relationships, enforcing data integrity and efficiently executing complex queries and update operations through Structured Query Language (SQL) instructions.

For the implementation of GeoMedia, Microsoft SQL Server (MSSQL) has been selected as the underlying relational database management system due to its robustness, scalability, and comprehensive support for enterprise-level data management. In addition to standard relational database capabilities, SQL Server provides advanced mechanisms for transaction management, data indexing and performance optimization.


MSSQL is freely available through the Express edition, which introduces some limitations in terms of storage (up to 10 GB per database), performance, and administration tooling. However, for the scope of this project, these limitations are considered acceptable. If the platform needs to scale in the future, the GeoMedia database can either be split across multiple instances, as discussed in the microservices section (Figure \ref{fig:UML-arch-micro}), or the system can be migrated to the Standard edition by acquiring the appropriate license.

\section{Database}

The database designed for the application follows some well known concepts such as

The database schema designed for the GeoMedia framework follows established relational database design principles, with particular attention to normalization, data integrity, and the appropriate definition of relationships between entities.
Specifically, the main principles adopted during the design are:

\begin{itemize}

    \item \textbf{Normalization}: the database schema is structured to minimize data redundancy and prevent insertion, update, and deletion anomalies, while maintaining clear relationships between entities.

    \item \textbf{Referential integrity}: primary and foreign key constraints are used to ensure that relationships between related entities remain consistent and that references to non-existent records are avoided.

    \item \textbf{ACID properties}: database transactions follow the principles of Atomicity, Consistency, Isolation, and Durability, ensuring that operations are executed reliably and that the database remains in a valid state.

    \item \textbf{Data integrity and consistency}: constraints such as \texttt{NOT NULL}, \texttt{UNIQUE}, and \texttt{CHECK} are employed where appropriate to prevent invalid or inconsistent data from being stored.

\end{itemize}

Moreover, the database schema was designed to facilitate the extensibility of the system by allowing additional entities and relationships to be introduced without requiring substantial modifications to the existing structure.

The Entity–Relationship (ER) diagram in Figure~\ref{fig:geomedia_ER} provides an overview of the main entities and their relationships within the GeoMedia database. For readability, entity attributes are omitted from the diagram, as several entities contain a considerable number of attributes.
%  The detailed attributes and corresponding database schema are instead defined at the implementation level.

\begin{figure}[H]
    \centering
    \includesvg[width=0.75\textwidth]{images/GeoMedia_ER.svg}
    \caption{GeoMedia Entity-Relationship diagram}
    \label{fig:geomedia_ER}
\end{figure}

The entities individuated are:
\begin{itemize}
    \item \textbf{Users}: which stores respective user informations and login credentials such as password and email. Moreover, create a reference to posts, reactions and eventual post reports.
    \item \textbf{Posts}: storing post information and then representing the core entity for integration and references.
    \item \textbf{Collections}: that works as containers for posts, dictating logic properties on posts and allowing special features (e.g. sequential posts).
    \item \textbf{Hypermedia}: operates as register with path of hypermedia stored by server layer in the file system.
    \item \textbf{Reactions}: actually implements just the `like' reaction, extendable to comments and others interaction among users and posts.
    \item \textbf{Reports}: store the report raised by users referenced to posts and then enable the moderation process of the platform.
\end{itemize}


In order to ensure the separation of GeoMedia's data from other databases hosted on the same database server, a dedicated database is created for the application using the \texttt{CREATE DATABASE GEOMEDIA} instruction. The database contains all the tables required to store the application's data, together with the corresponding stored procedures and functions. 
Through stored procedures (precompiled SQL routines), data can be manipulated through more complex operations, while business logic can be encapsulated within the database. This approach improves the maintainability and extensibility of the system, while potentially improving performance by reducing the amount of data processing required at the application level. Finally, the stored procedures provide the required results when requested by the application server as an example shown in Listing where the .





Instead for example of database structure, the SQL definition of the \textit{Posts} table is shown in Listing~\ref{lst:db_post_table}.


\lstset{
    language=SQL,
    basicstyle=\ttfamily\small,
    keywordstyle=\color{red},
    commentstyle=\color{green!60!black},
    stringstyle=\color{red},
    showstringspaces=false,
    frame=single,
    breaklines=true
}


\begin{lstlisting}[
    caption={SQL definition of the Posts table.},
    label={lst:db_post_table}
]

CREATE TABLE GEOMEDIA.dbo.POSTS (
	ID int IDENTITY(1, 1) NOT NULL,
	AUTHOR_ID int NOT NULL, -- author of post
	TITLE varchar(50) NOT NULL, --title of the post
	COMMENT varchar(MAX) NULL,
	DC datetime NULL,
	DM datetime DEFAULT getdate() NULL,
	LATITUDE float NULL, --latitude coordinate
	LONGITUDE float NULL, --longitude coordinate
	ALTITUDE float NULL, -- altitude for future developments
	VISIBILITY_AREA_KM float, -- visibility range
	EXCL_DATE_START datetime NULL, --exclusivity time start
	EXCL_DATE_END datetime NULL, --exclusivity time end
	RECURRENT int NULL, -- recurrency (0 = no reccurent, 1 = monthly, 2 = yearly)
	COLLECTION_ID int, 
	VIEWERS nvarchar(MAX) NULL, -- exclusivity viewers
	CONSTRAINT PK__POSTS__3214EC27ECDD25A2 PRIMARY KEY (ID),
    CONSTRAINT FK__POSTS__AUTHOR_ID__4F7CD00D FOREIGN KEY (AUTHOR_ID) REFERENCES GEOMEDIA.dbo.USERS(ID);
    CONSTRAINT FK__POSTS__AUTHOR_ID__7C4F7684 FOREIGN KEY (AUTHOR_ID) REFERENCES GEOMEDIA.dbo.USERS(ID);
    CONSTRAINT FK__POSTS__COLLECTIO__7D439ABD FOREIGN KEY (COLLECTION_ID) REFERENCES GEOMEDIA.dbo.COLLECTIONS(ID);
);

\end{lstlisting}


\section{Sequential post implementation}

The algorithm responsible for managing post sequentiality is implemented in the stored procedure \texttt{POST\_GET\_MAP}. The procedure is responsible for retrieving the posts available to the requesting user, applying the visibility and temporal constraints defined by the application and finally determining the appropriate order of sequential posts. Meanwhile, the radius-based area filter is performed by the backend, as described in Section~\ref{haversine_formula}, on the dataset returned by the stored procedure. \\

In particular, the procedure handles sequential posts by identifying the first post in the defined sequence that has not yet been viewed by the requesting user, allowing the framework to progressively expose posts according the sequence defined in the respective collections selected by user. \\
The sequentiality logic is implemented through a CTE (Common Table Expression), that extracts the ordering information associated with collections from \texttt{SEQUENTIALS} field stored in JSON format, then applying the SQL Server built-in function \texttt{OPENJSON} parse the structure in form of table with corresponding  \texttt{order\_id} and \texttt{post\_id} values. \\
The orders values are stored temporarily in the \texttt{@SEQQ} table variable, subsequently used to associate each posts with its position within the sequence. INstead, a second CTE \texttt{NEXTPOST} identifies the first post in the sequence that has not yet been viewed by the requesting users. \\
This is achieved by comparing the posts associated with the collections against the records stored in the \texttt{INTERACTIONS} table, where the value \texttt{VIEW} in field \texttt{SCOPE} marks the post as seen by the user. \\
The results are ordered according to the sequential \texttt{order\_id}, ensuring that the next available post is selected according to the sequence defined by collection owner.
Finally, the procedure combines the next post to be displayed of respective collections with sequentiality enabled, and posts that are not associated with a sequence. Additional filters are applied to respect account visibility (both post than collection), temporal restrictions and respective recurrence rules.

The implementation of this logic is shown in Listing~\ref{lst:db_sp_post_seq}.


\lstset{
    language=SQL,
    basicstyle=\ttfamily\footnotesize, %\small?
    keywordstyle=\color{red},
    commentstyle=\color{green!60!black},
    stringstyle=\color{red},
    showstringspaces=false,
    frame=single,
    breaklines=true
}

\begin{lstlisting}[
    caption={SQL implementation for post sequentiality.},
    label={lst:db_sp_post_seq}
]
CREATE PROCEDURE [dbo].[POST_GET_MAP]
@UID INT, @COLLECTIONS_CHOSEN NVARCHAR(MAX)= NULL
AS
BEGIN
    DECLARE @SEQQ TABLE (order_id int,post_id int);

    ;WITH SEQ AS(
        SELECT j.order_id, j.post_id
        FROM COLLECTIONS c
        CROSS APPLY OPENJSON(c.SEQUENTIALS, '$')
        WITH (
            order_id INT,
            post_id INT
        ) j
        WHERE c.ID IN (SELECT VALUE FROM OPENJSON(@COLLECTIONS_CHOSEN, '$'))
    )
    INSERT INTO @SEQQ(order_id, post_id) 
    SELECT order_id, post_id FROM SEQ;
    --- SEQUENTIALITY POSTS
    ;
    WITH NEXTPOST AS(--- FIRST POST NOT VIEWED
        SELECT TOP 1 P.*, SEQ.order_id, SEQ.post_id
        FROM POSTS P
        LEFT JOIN @SEQQ SEQ ON SEQ.post_id = P.ID
        WHERE P.COLLECTION_ID IN (SELECT VALUE FROM OPENJSON(@COLLECTIONS_CHOSEN, '$'))
        AND P.ID NOT IN (
            SELECT POST_ID FROM INTERACTIONS WHERE OWNERID = @UID AND SCOPE = 'VIEW'
        )
        ORDER BY SEQ.order_id ASC
    )
    SELECT P.ID, P.AUTHOR_ID, P.TITLE, P.COMMENT, P.DC, P.DM, P.LATITUDE, P.LONGITUDE, P.VISIBILITY_AREA_KM, P.EXCL_DATE_START, P.EXCL_DATE_END, P.RECURRENT, P.COLLECTION_ID, P.VIEWERS, SEQ.order_id, SEQ.post_id, C.COLOR, C.ICON, C.TITLE AS CATEGORY_NAME
    FROM POSTS P
    LEFT JOIN @SEQQ SEQ ON SEQ.post_id = P.ID
    INNER JOIN COLLECTIONS C ON C.ID = P.COLLECTION_ID
    WHERE P.COLLECTION_ID IN ( SELECT VALUE FROM OPENJSON(@COLLECTIONS_CHOSEN, '$'))
    AND P.ID IN ( SELECT POST_ID FROM INTERACTIONS WHERE OWNERID = @UID AND SCOPE = 'VIEW' )
    ---- VIEWED POST ^
    UNION
        SELECT P.ID, P.AUTHOR_ID, P.TITLE, P.COMMENT, P.DC, P.DM, P.LATITUDE, P.LONGITUDE, P.VISIBILITY_AREA_KM, P.EXCL_DATE_START, P.EXCL_DATE_END, P.RECURRENT, P.COLLECTION_ID, P.VIEWERS, P.order_id, P.post_id, C.COLOR, C.ICON, C.TITLE AS CATEGORY_NAME
        FROM NEXTPOST P
    INNER JOIN COLLECTIONS C ON C.ID = P.COLLECTION_ID
    UNION
    --- UNION WITH POST WITHOUT SEQUENCE
    SELECT P.ID, P.AUTHOR_ID, P.TITLE, P.COMMENT, P.DC, P.DM, P.LATITUDE, P.LONGITUDE, P.VISIBILITY_AREA_KM, P.EXCL_DATE_START, P.EXCL_DATE_END, P.RECURRENT, P.COLLECTION_ID, P.VIEWERS, NULL AS order_id, NULL AS post_id, C.COLOR, C.ICON, C.TITLE AS CATEGORY_NAME
    FROM POSTS P
    INNER JOIN COLLECTIONS C ON C.ID = P.COLLECTION_ID
    AND (
        --COLLECTION TIME LIMITING CHECK
            (C.RECURRENT IS NULL and (GETDATE()> ISNULL(C.EXCL_DATE_START, DATEADD(DAY,-1, GETDATE()))
                OR GETDATE() < ISNULL(C.EXCL_DATE_END, DATEADD(DAY, 1, GETDATE())))
            )
        -- NO C.RECURRENTCY
            OR  (C.RECURRENT = 1 AND (DAY(GETDATE()) BETWEEN DAY(C.EXCL_DATE_START) AND DAY(C.EXCL_DATE_START))) --EACH DAYS, IGNORING MONTHS
                OR (C.RECURRENT = 2 AND  (DATEFROMPARTS(2000, MONTH(GETDATE()), DAY(GETDATE())) BETWEEN DATEFROMPARTS(2000, MONTH(C.EXCL_DATE_START), DAY(C.EXCL_DATE_START))
                    AND DATEFROMPARTS(2000, MONTH(C.EXCL_DATE_END), DAY(C.EXCL_DATE_END)))
            ) -- EVERY YEAR, SAME DAYS SAME MONTH
        )
    WHERE C.SEQUENTIALS IS NULL -- NO SEQUENTIALITY INDICATED (GENERAL  POSTS)
        AND ( P.VIEWERS IS NULL OR P.VIEWERS = '[]' OR @UID IN (
            SELECT VALUE FROM OPENJSON(P.VIEWERS, '$'))
            AND
            (
                (P.RECURRENT IS NULL
                    and (GETDATE()> ISNULL(P.EXCL_DATE_START, DATEADD(DAY,-1, GETDATE()))
                        OR GETDATE() < ISNULL(P.EXCL_DATE_END, DATEADD(DAY, 1, GETDATE())))
                    --IF NULL, MAKE IT FROM YESTERDAY TILL TOMORROW
                )
                -- NO P.RECURRENTCY
                    OR 
                (P.RECURRENT = 1 AND (DAY(GETDATE()) BETWEEN DAY(P.EXCL_DATE_START) AND DAY(P.EXCL_DATE_START))) --EACH DAYS, IGNORING MONTHS
                    OR (P.RECURRENT = 2 AND 
                    (DATEFROMPARTS(2000, MONTH(GETDATE()), DAY(GETDATE())) BETWEEN DATEFROMPARTS(2000, MONTH(P.EXCL_DATE_START), DAY(P.EXCL_DATE_START))
                    AND DATEFROMPARTS(2000, MONTH(P.EXCL_DATE_END), DAY(P.EXCL_DATE_END)))
                )
                -- EVERY YEAR, SAME DAYS SAME MONTH
            ) AND ( @UID IN (
                    --CHECK IF THE USER REQUESTING IS ENABLED TO THAT CATEGORY
                    SELECT VALUE FROM OPENJSON(C.VIEWERS, '$')
                ) OR C.VIEWERS = '[]'
            )
        )
        AND C.ID IN ( SELECT VALUE FROM OPENJSON(@COLLECTIONS_CHOSEN, '$'))
        -- NB: THE VISIBILITY_AREA_KM FILTER IS MADE ON BACKEND OF THE SERVER
END;
\end{lstlisting}


\subsection{Docker Deployment}

